%%
%% This is file `sample-acmsmall-submission.tex',
%% generated with the docstrip utility.
%%
%% The original source files were:
%%
%% samples.dtx  (with options: `all,journal,bibtex,acmsmall-submission')
%% 
%% IMPORTANT NOTICE:
%% 
%% For the copyright see the source file.
%% 
%% Any modified versions of this file must be renamed
%% with new filenames distinct from sample-acmsmall-submission.tex.
%% 
%% For distribution of the original source see the terms
%% for copying and modification in the file samples.dtx.
%% 
%% This generated file may be distributed as long as the
%% original source files, as listed above, are part of the
%% same distribution. (The sources need not necessarily be
%% in the same archive or directory.)
%%
%%
%% Commands for TeXCount
%TC:macro \cite [option:text,text]
%TC:macro \citep [option:text,text]
%TC:macro \citet [option:text,text]
%TC:envir table 0 1
%TC:envir table* 0 1
%TC:envir tabular [ignore] word
%TC:envir displaymath 0 word
%TC:envir math 0 word
%TC:envir comment 0 0
%%
%% The first command in your LaTeX source must be the \documentclass
%% command.
%%
%% For submission and review of your manuscript please change the
%% command to \documentclass[manuscript, screen, review]{acmart}.
%%
%% When submitting camera ready or to TAPS, please change the command
%% to \documentclass[sigconf]{acmart} or whichever template is required
%% for your publication.
%%
%%
\documentclass[acmsmall,screen,review]{acmart}
%%
%% \BibTeX command to typeset BibTeX logo in the docs
\AtBeginDocument{%
  \providecommand\BibTeX{{%
    Bib\TeX}}}

%% Rights management information.  This information is sent to you
%% when you complete the rights form.  These commands have SAMPLE
%% values in them; it is your responsibility as an author to replace
%% the commands and values with those provided to you when you
%% complete the rights form.
\setcopyright{acmlicensed}
\copyrightyear{2026}
\acmYear{2026}
\acmDOI{XXXXXXX.XXXXXXX}

%%
%% These commands are for a JOURNAL article.
\acmJournal{JACM}
\acmVolume{37}
\acmNumber{4}
\acmArticle{111}
\acmMonth{8}

%%
%% Submission ID.
%% Use this when submitting an article to a sponsored event. You'll
%% receive a unique submission ID from the organizers
%% of the event, and this ID should be used as the parameter to this command.
%%\acmSubmissionID{123-A56-BU3}

%%
%% For managing citations, it is recommended to use bibliography
%% files in BibTeX format.
%%
%% You can then either use BibTeX with the ACM-Reference-Format style,
%% or BibLaTeX with the acmnumeric or acmauthoryear sytles, that include
%% support for advanced citation of software artefact from the
%% biblatex-software package, also separately available on CTAN.
%%
%% Look at the sample-*-biblatex.tex files for templates showcasing
%% the biblatex styles.
%%

%%
%% The majority of ACM publications use numbered citations and
%% references.  The command \citestyle{authoryear} switches to the
%% "author year" style.
%%
%% If you are preparing content for an event
%% sponsored by ACM SIGGRAPH, you must use the "author year" style of
%% citations and references.
%% Uncommenting
%% the next command will enable that style.
%%\citestyle{acmauthoryear}


\begin{document}

%%
%% The "title" command has an optional parameter,
%% allowing the author to define a "short title" to be used in page headers.
\title{Verifying Distributed Algorithms with Interactive Theorem Provers: A Survey of Approaches}

%%
%% The "author" command and its associated commands are used to define
%% the authors and their affiliations.
%% Of note is the shared affiliation of the first two authors, and the
%% "authornote" and "authornotemark" commands
%% used to denote shared contribution to the research.
\author{Maxwell Moir}
\affiliation{%
  \institution{University of Canterbury}
  \city{Christchurch}
  \country{New Zealand}}
\email{mmo199@uclive.ac.nz}


%%
%% By default, the full list of authors will be used in the page
%% headers. Often, this list is too long, and will overlap
%% other information printed in the page headers. This command allows
%% the author to define a more concise list
%% of authors' names for this purpose.
\renewcommand{\shortauthors}{Moir}

%%
%% The abstract is a short summary of the work to be presented in the
%% article.
\begin{abstract}

  Formal verification is a topical field in modern Computer Science. 
  Proving the correctness of programs is not only elegant in a theoretical sense, but is often necessary in mission-critical applications.
  Due to the difficulty in formalizing programs, reliability in industry is often left to automated testing. 
  However, in specific circumstances where correctness must be ensured, formal proofs may be required. 
  This survey reviews the different approaches to verifying the correctness of distributed algorithms.

\end{abstract}


%%
%% The code below is generated by the tool at http://dl.acm.org/ccs.cfm.
%% Please copy and paste the code instead of the example below.
%%
\begin{CCSXML}
<ccs2012>
   <concept>
       <concept_id>10003752.10003790.10002990</concept_id>
       <concept_desc>Theory of computation~Logic and verification</concept_desc>
       <concept_significance>500</concept_significance>
       </concept>
 </ccs2012>
\end{CCSXML}

\ccsdesc[500]{Theory of computation~Logic and verification}


%%
%% Keywords. The author(s) should pick words that accurately describe
%% the work being presented. Separate the keywords with commas.
\keywords{distributed algorithms}

%%
%% This command processes the author and affiliation and title
%% information and builds the first part of the formatted document.
\maketitle

\section{Introduction}

With increasing societal dependance on software, ensuring the effective operation of the tools we use has become more important.

% Motivations
%   Why its hard, why it matters, real world uses
%   Difficulties - Distributed systems are notorious for being difficult to verify. 
%   There are many edge cases and etc... that make these systems
One method for reasoning about the correctness of programs is formal verification.
Distributed systems underpin many of the applications that we use every day.
Distributed systems are notoriously difficult to verify due to ...

% Rise of proof assistant verification
% Recently, there has been a rise in verification with Interactive Theorem Provers.

% What this survey contributes:
%   Overview of verification approaches
%   comparison 
%   challenges and open problems
This survey aims to explore attempts to formally verify distributed algorithms, both theoretically and in practice.
It provides a taxonomy of different theorem-based verfication approaches and frameworks, reviews several common tools, and contrasts with other verification approaches.

% Scope
Other approaches to the verification of distributed algorithms are outside the scope of this survey, except when used as comparison. 


\section{Preliminary}

\subsection{Distributed Protocols}
Distributed computing is ...

\subsection{Interactive Theorem Provers}

\subsection{Verification}




\section{Approaches}

Nam id fermentum dui. Suspendisse sagittis tortor a nulla mollis, in
pulvinar ex pretium. Sed interdum orci quis metus euismod, et sagittis
enim maximus. Vestibulum gravida massa ut felis suscipit
congue. Quisque mattis elit a risus ultrices commodo venenatis eget
dui. Etiam sagittis eleifend elementum.

There are several common frameworks used in the verification of distributed protocols. 
Each of which has its own benefits and detractions.


Common Frameworks:
Verdi \cite{wilcoxVerdiFrameworkImplementing2015}
Iron Fleet \cite{hawblitzelIronFleetProvingPractical2015}
Diesel \cite{sergeyProgrammingProvingDistributed2017}
Veil \cite{pirleaVeilFrameworkAutomated2025a}
TLAPS \cite{cousineau2012tla}

Common Languages:
Coq
Lean
Isabelle
Dafny

Common Problems
Paxos
Raft

Verified system examples

Byzantine fault tolerance

Compositional and modular verification


Proof techniques:
Invariant
refinement

Adjacents


\section{Applications}
Nam id fermentum dui. Suspendisse sagittis tortor a nulla , in
pulvinar ex pretium. Sed interdum orci quis metus euismod, et sagittis

\subsection{Verified protocols}

Nam id fermentum dui. Suspendisse sagittis tortor a nulla , in
pulvinar ex pretium. Sed interdum orci quis metus euismod, et sagittis
enim maximus. Vestibulum gravida massa ut felis suscipit
congue. Quisque mattis elit a risus ultrices commodo venenatis eget
dui. Etiam sagittis eleifend elementum.

\subsection{Verified systems}

Nam id fermentum dui. Suspendisse sagittis tortor a nulla , in
pulvinar ex pretium. Sed interdum orci quis metus euismod, et sagittis
enim maximus. Vestibulum gravida massa ut felis suscipit
congue. Quisque mattis elit a risus ultrices commodo venenatis eget
dui. Etiam sagittis eleifend elementum.

r plants. Non-safety-critical applications include a web browser and few social applications: a conferencing system, a social media platform.
In spite of the potential of today’s verification technology and its successes as documented here, formal verification is still not mainstream. The key obstacle is the high entry level: each project needs one or more verification experts. Companies may feel that it will be easier to find many testers than a few verification-savvy developers.
Another concern is the rate of change of modern software systems. It is now common to practice “continuous development”, producing new versions once or more daily. In principle every such version requires re-verification. The good news, reported by several projects (seL4, Cocon, Amazon s2n, iFACTS, and Dutch tunnel) is that re-verification costs much less than the original verification. seL4 reports that verifying a new and large feature (a new data structure for API calls) cost less than 10% of the initial effort. The re-verification of small changes in iFACTS was completed in seconds, while its initial verification run took three hours.
An alternative to re-verification is reuse of verified elements. Reviewing the projects revealed a number of cases of successful reuse, for example in OS projects (mCertiKos 6.5.10, Hyper-V, and PikeOS). Particularly interesting is the reuse by VeriCert of
\section{Alternative Verification Approaches}

Model checking, 


\section{Future Directions}

Nam id fermentum dui. Suspendisse sagittis tortor a nulla , in
pulvinar ex pretium. Sed interdum orci quis metus euismod, et sagittis
enim maximus. Vestibulum gravida massa ut felis suscipit
congue. Quisque mattis elit a risus ultrices commodo venenatis eget
dui. Etiam sagittis eleifend elementum.


\section{Conclusion}

Nam id fermentum dui.  sagittis tortor a nulla mollis, in
pulvinar ex pretium. Sed interdum orci quis metus euismod, et sagittis
enim maximus. Vestibulum gravida massa ut felis suscipit
congue. Quisque mattis elit a risus ultrices commodo venenatis eget
dui. Etiam sagittis eleifend elementum. Test 01

\section{Methodology}

\subsection{Literature Search Strategy}

To identify relevant papers, we used a systematic literature search strategy. 
We chose to query two databases, the ACM digital library and IEEE Xplore. 
These databases were accessed on April 17 2026.

To ensure the quality of of this review, we decided to follow the PRISMA framework to inform our search strategy.
The search strategy can be summarized into multiple steps.
First, we curated relevant keywords and constructed seach queries to compile relevant papers based on the guiding questions of the survey.
This led to the following seach query:

("proof assistant" OR "theorem prover" OR "mechanized verification" OR "formal verification") AND ("distributed protocol" OR "distributed algorithm")

This query yeilded 213 results from the ACM Digital Library and 49 results from IEEE Xplore.

\subsection{Inclusion and Exclusion Criteria}
These results were then screened based on standard inclusion and exclusion requirements.
This reduced the result count to 37 for ACM and XX for IEEE Xplore. 
From this point, the results were expanded by adding commonly cited papers.



%% The next two lines define the bibliography style to be used, and
%% the bibliography file.
\bibliographystyle{ACM-Reference-Format}
\bibliography{sample-base}

\end{document}
\endinput

